% T037: exact side-length model, A=(0,0), B=(c,0).
% Arguments a=BC,b=CA,c=AB,scale,mode,BC label,CA label,AB label,caption.
% Modes 0 area, 1 common altitude, 2 tangent pairs, 3 angle, 4 altitude example, 5 area pieces.
\providecommand{\THeronFigure}[9]{%
\begin{center}
\begin{tikzpicture}[scale=#4,font=\small,line join=round,line cap=round]
\pgfmathsetmacro{\tq}{((#2)^2+(#3)^2-(#1)^2)/(2*(#3))}
\pgfmathsetmacro{\th}{sqrt((#2)^2-\tq^2)}
\pgfmathsetmacro{\tp}{((#1)+(#2)+(#3))/2}
\pgfmathsetmacro{\tr}{(#3)*\th/(2*\tp)}
\pgfmathsetmacro{\tix}{((#2)*(#3)+(#3)*\tq)/(2*\tp)}
\coordinate (A) at (0,0);\coordinate (B) at (#3,0);
\coordinate (C) at (\tq,\th);\coordinate (H) at (\tq,0);
\coordinate (I) at (\tix,\tr);
\coordinate (D) at ($(B)!(I)!(C)$);
\coordinate (E) at ($(C)!(I)!(A)$);
\coordinate (F) at ($(A)!(I)!(B)$);
\ifnum#5=0
 \fill[blue!10] (A)--(B)--(C)--cycle;
 \node at ({(#3+\tq)/3},{\th/3}) {$S$};
\fi
\ifnum#5=7
 \fill[blue!10] (A)--(B)--(C)--cycle;
 \pic[draw,angle radius=3mm] {right angle=A--C--B};
\fi
\ifnum#5=1
 \fill[blue!13] (A)--(H)--(C)--cycle;
 \fill[orange!18] (B)--(H)--(C)--cycle;
\fi
\ifnum#5=5
 \fill[blue!13] (I)--(B)--(C)--cycle;
 \fill[orange!18] (I)--(C)--(A)--cycle;
 \fill[green!15!white] (I)--(A)--(B)--cycle;
\fi
\draw[thick] (A)--(B)--(C)--cycle;
\ifnum#5=2\else
 \path (B)--node[pos=.82,right=5pt] {$#6$} (C);
 \path (C)--node[pos=.25,left=5pt] {$#7$} (A);
 \ifnum#5=1\else
  \ifnum#5=4
   \path (A)--node[midway,below=28pt] {$#8$} (B);
  \else
   \ifnum#5=5
    \path (A)--node[midway,below=28pt] {$#8$} (B);
   \else
    \path (A)--node[midway,below=12pt] {$#8$} (B);
   \fi
  \fi
 \fi
\fi
\foreach \v/\where in {A/below left,B/below right,C/above}
 {\fill (\v) circle[radius=1.4pt];\node[\where=4pt] at (\v) {$\v$};}
\ifnum#5=1
 \draw[very thick,CProof] (C)--node[midway,right=4pt] {$h$} (H);
 \pic[draw,angle radius=2.8mm] {right angle=C--H--B};
 \fill (H) circle[radius=1.2pt];\node[below=4pt] at (H) {$D$};
 \path (A)--node[midway,below=16pt] {$x$} (H);
 \path (H)--node[midway,below=16pt] {$c-x$} (B);
\fi
\ifnum#5=4
 \draw[very thick,CProof] (C)--node[midway,right=4pt] {$h$} (H);
 \pic[draw,angle radius=2.8mm] {right angle=C--H--B};
 \fill (H) circle[radius=1.2pt];\node[below=4pt] at (H) {$D$};
\fi
\ifnum#5=3
 \coordinate (K) at ($(B)!(A)!(C)$);
 \draw[CProof,dashed,thick] (A)--node[pos=.55,above=6pt] {$h_a$} (K);
 \pic[draw,angle radius=2.8mm] {right angle=A--K--C};
 \pic[draw,CProof,angle radius=9mm] {angle=A--C--B};
\fi
\ifnum#5=2\THeronIncircleDetails{2}\fi
\ifnum#5=5\THeronIncircleDetails{5}\fi
\ifnum#5=0
 \node[anchor=north,align=center,text width=11cm] at (#3/2,{-1.25/#4}) {#9};
\else
 \node[anchor=north,align=center,text width=11cm] at (#3/2,{-1.65/#4}) {#9};
\fi
\end{tikzpicture}
\end{center}}

\providecommand{\THeronIncircleDetails}[1]{%
\draw[CStatement,thick] (I) circle[radius=\tr];
\draw[gray,dashed] (I)--(A) (I)--(B) (I)--(C);
\draw[very thick,CStatement] (I)--(D) (I)--(E) (I)--node[midway,right=3pt] {$r$} (F);
\pic[draw,angle radius=2mm] {right angle=I--D--C};
\pic[draw,angle radius=2mm] {right angle=I--E--A};
\pic[draw,angle radius=2mm] {right angle=I--F--B};
\foreach \v/\where in {D/above right,E/above left,F/below}
 {\fill (\v) circle[radius=1.3pt];\node[\where=4pt] at (\v) {$\v$};}
\fill (I) circle[radius=1.6pt];\node at ($(I)+(1.3,0)$) {$I$};
\ifnum#1=2
 \draw[very thick,CStatement] (A)--node[midway,left=6pt] {$x$} (E)
  (A)--node[midway,below=16pt] {$x$} (F);
 \draw[very thick,CExample] (B)--node[midway,right=6pt] {$y$} (D)
  (B)--node[midway,below=16pt] {$y$} (F);
 \draw[very thick,CProof] (C)--node[midway,left=6pt] {$z$} (E)
  (C)--node[midway,right=6pt] {$z$} (D);
\fi}

\begin{statementbox}
  \textbf{\textcolor{CStatement}{这个结论解决什么问题？}}

  题目只告诉你三角形的\textbf{三条边长}，却没有告诉你高、角度或坐标，怎样直接求出面积？
  海伦公式把原本需要求高或角度的题，变成\textbf{加、减、乘、开平方}。

\THeronFigure{13}{14}{15}{0.58}{0}{a}{b}{c}{三边确定，面积 $S$ 就确定；不必把未知的高或角先算出来。}
  \medskip
  \textbf{\textcolor{CStatement}{条件与字母含义}}
  \begin{description}[style=nextline,leftmargin=2.8em,labelsep=0.5em,itemsep=0.35em]
    \item[\textbf{条件一：}] $ABC$ 是一个真正的三角形；$a=BC>0$，$b=CA>0$，$c=AB>0$。
    \item[\textbf{条件二：}] $a+b>c$，$b+c>a$，$c+a>b$（任意两边之和大于第三边）。
    \item[\textbf{记号：}] 面积记为 $S>0$，半周长记为 $p=\dfrac{a+b+c}{2}$。
  \end{description}

  \textbf{\textcolor{CStatement}{结论一：已知三边，直接求面积（核心）}}
  \[
    \boxed{\displaystyle S=\sqrt{p(p-a)(p-b)(p-c)},\qquad p=\frac{a+b+c}{2}.}
  \]
  四个量依次是\textbf{半周长}、\textbf{半周长减第一边}、\textbf{减第二边}、\textbf{减第三边}。

  \textbf{\textcolor{CStatement}{结论二：不设半周长的等价写法}}
  \[
   \boxed{\displaystyle 16S^2=(a+b+c)(a+b-c)(a-b+c)(-a+b+c).}
  \]
  \textbf{用途：}在推导中看见 ``两个平方相减'' 时，上面这四个因式会自然出现；它们不是需要另外死记的一套公式。

\par\noindent\begin{minipage}{\linewidth}
  \textbf{\textcolor{CStatement}{结论三：求内切圆半径（延伸）}}
  设内切圆半径为 $r$，连接内心 $I$ 与三个顶点：
\THeronFigure{13}{14}{15}{0.44}{5}{a}{b}{c}{三块面积相加：$S=ar/2+br/2+cr/2=pr$。}
\end{minipage}\par
  因为 $S=pr$，所以
  \[
    \boxed{\displaystyle r=\frac{S}{p}=\sqrt{\frac{(p-a)(p-b)(p-c)}{p}}.}
  \]

  \textbf{\textcolor{CStatement}{最重要的适用范围：}}
  \textbf{三边都能确定而高、角不方便求}时优先考虑；\textbf{直角三角形已经有两条直角边}时，直接用 ``底$\times$高$\div 2$'' 往往更快。
\end{statementbox}
